
\subsection{Contents}

This book includes selected articles by John Evjen related to the works of Georg Sverdrup and Sven Oftedal. Evjen is a witness to what is now the Fundamental Principles of the Association of Free Lutheran Congregations (AFLC) with an emphasis on Principle \#1. A copy of the AFLC's Fundamental Principles is included in the front matter.

This is not a biography of Evjen, Sverdrup, or Oftedal. Nor is it a history of the AFLC. However, Evjen's works from 1910 to 1918 are still with us today. Together, the five sections ask if the AFLC principles still serve as the foundation for living congregations or if ``spirit and life'' has become a slogan.   

This book contains five sections:

\begin{itemize}

\item Section 1: \textit{A Layman's Introduction} by Aaron Peder Dahlen.

\item Section 2: \textit{A Contribution to the Understanding of Professor Georg Sverdrup and His Selected Writings}: The text originally appeared in \textit{Folkebladet} on May 4, 11, and 18, 1910. It was later bound and published as a small pamphlet by the Free Church Book Concern of Minneapolis, Minnesota. 

This section is important as it captures Evjen's optimism. It was written immediately after the release of the first two [of six] volumes of Professor Sverdrup's \textit{Sketches and Addresses} (\textit{Skisser og Afhandlinger}) edited by Andreas Helland. Evjen's pamphlet reads like an overly enthusiastic review of the twin volumes. We can feel Evjen's relief. He is settling into the joy that Sverdrup's words are now accessible. He projects that it is only a matter of time before the Lutheran Free Church (FLC) will read and understand Sverdrup's writings. We know better.  

\item Section 3: \textit{Guidance: The Principles of the Lutheran Free Church}: This is a collection of writings by Sverdrup and Oftedal compiled by Evjen. It was published in book form in 1914 by the Free Church Book Concern of Minneapolis, Minnesota.

\item Section 4: \textit{The Position}: This is an 11-part argument written by Evjen. It appears in \textit{Folkebladet} from March 13 to May 22, 1918.

\item Section 5: \textit{The Editors}: This includes materials written by the \textit{Folkebladet's} editors. Their writings are important to understanding Evjen's exchange with The Editors in ``The Position.'' 
  
\end{itemize}

Note that \textit{Guidance} contains no commentary by Evjen---not even a foreword explaining why these particular writings were selected. The inclusion of Oftedal's Hennepin County affidavit is also puzzling. From ``The Position,'' we learn that Evjen used \textit{Guidance} as a textbook while teaching at seminary.

